\documentclass[12pt]{article}
\usepackage[utf8]{inputenc}
\usepackage[T1]{fontenc}
\usepackage{hyperref}
\title{Yabarana Language Resources for K-12 Curriculum Development}
\author{Compiled by Claude Sonnet 5.5 (claude-sonnet-5-5) for Andr\'as Kornai}
\date{}
\begin{document}
\maketitle

\section{Introduction}
Yabarana (also Yawarana) is a Cariban language of the Manapiare River basin in Amazonas, Venezuela. Estimates of its speakers vary widely: from fewer than 25 fluent speakers to about 150. An AILLA record puts fluent first-language speakers at about 20--30, mostly over 60. Sources also disagree on whether it is a separate language or a variety close to Mapoyo. Very little material exists, and this report lists only what the searches turned up.

\section{Primary Educational Resources}
The searches found no primers, textbooks or school materials in Yabarana. Nothing found indicates that the language is taught in schools.

\section{Linguistic and Reference Materials}
\begin{itemize}
\item Wikipedia overview of the language: \url{https://en.wikipedia.org/wiki/Yabarana_language}
\item Glottolog entry (yaba1248): \url{https://glottolog.org/resource/languoid/id/yaba1248.bigmap.html}
\item ASJP word list entry: \url{https://asjp.clld.org/languages/YABARANA}
\item Endangered Languages Catalogue entry: \url{https://endangeredlanguages.com/lang/1649?hl=pt-br}
\item Spike Gildea's academic page, which lists his work on Cariban tense-aspect systems including Yabarana: \url{https://uoregon.academia.edu/SpikeGildea}
\end{itemize}
A search snippet says that before 2017 only three word lists and an undergraduate thesis on a few morphological points existed.

\section{Religious and Community Resources}
Joshua Project lists the Yabarana people in Venezuela: \url{https://joshuaproject.net/people_groups/15929/VE}. No scripture translation or community organisation was found.

\section{Digital and Audio Resources}
The Archive of the Indigenous Languages of Latin America (AILLA) holds a record that appears to contain Yabarana material (the search snippet's content was not checked): \url{https://islandora-ailla.lib.utexas.edu/islandora/object/ailla%3A119590}

\section{Curriculum Development Context}
With about 20 speakers, mostly elderly, a school curriculum would first need documentation: recordings, a practical orthography and a lexicon, made with the community and the AILLA holdings. Nothing found suggests an existing orthography standard or teaching programme. Any curriculum would be revitalisation-oriented, with Spanish as the language of schooling.

\end{document}